\let\negmedspace\undefined
\let\negthickspace\undefined
\documentclass[journal]{IEEEtran}
\usepackage[a5paper, margin=10mm, onecolumn]{geometry}
%\usepackage{lmodern} % Ensure lmodern is loaded for pdflatex
\usepackage{tfrupee} % Include tfrupee package
\setlength{\headheight}{1cm} % Set the height of the header box
\setlength{\headsep}{0mm}     % Set the distance between the header box and the top of the text
\usepackage{gvv-book}
\usepackage{gvv}
\usepackage{cite}
\usepackage{amsmath,amssymb,amsfonts,amsthm}
\usepackage{algorithmic}
\usepackage{graphicx}
\usepackage{textcomp}
\usepackage{xcolor}
\usepackage{txfonts}
\usepackage{listings}
\usepackage{enumitem}
\usepackage{mathtools}
\usepackage{gensymb}
\usepackage{comment}
\usepackage[breaklinks=true]{hyperref}
\usepackage{tkz-euclide} 
\usepackage{listings}
% \usepackage{gvv}                                        
\def\inputGnumericTable{}                                 
\usepackage[latin1]{inputenc}                                
\usepackage{color}                                            
\usepackage{array}                                            
\usepackage{longtable}                                       
\usepackage{calc}                                             
\usepackage{multirow}                                         
\usepackage{hhline}                                           
\usepackage{ifthen}                                           
\usepackage{lscape}
\usepackage{circuitikz}
\tikzstyle{block} = [rectangle, draw, fill=blue!20, 
    text width=4em, text centered, rounded corners, minimum height=3em]
\tikzstyle{sum} = [draw, fill=blue!10, circle, minimum size=1cm, node distance=1.5cm]
\tikzstyle{input} = [coordinate]
\tikzstyle{output} = [coordinate]
\begin{document}
\bibliographystyle{IEEEtran}
\vspace{3cm}
\title{5.3.11}
\author{AI25BTECH11017-BALU}
 \maketitle
% \newpage
% \bigskip
{\let\newpage\relax\maketitle}
\renewcommand{\thefigure}{\theenumi}
\renewcommand{\thetable}{\theenumi}
\setlength{\intextsep}{10pt} % Space between text and floats
\numberwithin{equation}{enumi}
\numberwithin{figure}{enumi}
\renewcommand{\thetable}{\theenumi}
\textbf{Question}:\\
Find $\vec{A}^2-5\vec{A}+6\vec{I}$ if 
\begin{align}
    \vec{A}=\myvec{2&&0&&1\\2&&1&&3\\1&&-1&&0}\
\end{align}
\\
\solution \\
Let us solve the given equation theoretically and then verify the solution computationally \\
According to the question, \\
We known that\\
\begin{align}
    \|\vec{A}-\lambda\vec{I}\|=0\\
    \left\|\myvec{2-\lambda&&0&&1\\2&&1-\lambda&&3\\1&&-1&&-\lambda}\right\|=0\\
    \lambda^3-3\lambda^2+4\lambda-3=0\
\end{align}
using Cayley-Hamilton theorem
\begin{align}
    \vec{A}^3-3\vec{A}^2+4\vec{A}-3\vec{I}=0\\
    \vec{A}^2-3\vec{A}+4\vec{A}=3\vec{A}^{-1}\\
    \vec{A}^2-5\vec{A}+6\vec{I}=3\vec{A}^{-1}-2\vec{A}+2\vec{I}=\myvec{1&&-1&&-3\\-1&&-1&&-10\\-5&&4&&4}
\end{align}
\end{document}